% !TeX root = ../../Book5.tex
% 纲要标题：## 附录 B　量子群、簇代数与范畴化
% 纲要位置：工作记录/计划纲要.md，第 4948 行。
% 纲要细目（保留原定写作范围）：
%
% * 量子群作为 Hopf 代数的例子；余乘法的一次修正、二维 Hecke 翻转算子及单位根模型的合成列与非分裂扩张
% * 簇变异、A2 编号周期与 Laurent 现象证明概要；以两次变异和差商解释新箭头及分母消去，区分通用互素性、特殊化和簇代数中的可逆性
% * Grothendieck 群、正合关系与分次 Euler 类；区分对象的类、实际扩张与函子，并说明 Hecke 规范基的分次解释范围
% * 简介不代替相应专题课程的完整理论
%
\chapter{量子群、簇代数与范畴化}
\label{b5:app:B}
\begin{chapterabstract}
	本附录选取量子 \(\mathfrak{sl}_2\)、五周期簇变异与 Grothendieck 群三个模型，分别讨论参数变形、有理坐标变换和正合列产生的环关系。
\end{chapterabstract}
\begin{learninggoals}
	\begin{itemize}
		\item 核验量子包络代数的关系与Hopf公式，判断普通张量翻转为何可能不等变。
		\item 计算二阶簇变量，区分单个算例的Laurent性与一般定理。
		\item 比较分裂 Grothendieck 群与按短正合列关系定义的 Grothendieck 群，并用张量函子解释乘法。
		\item 在单位根、参数退化及点代入中检查原有证明的假设。
	\end{itemize}
\end{learninggoals}

从两个代数独立的变量 \(x,y\) 出发，先令 \(x'=(1+y)/x\)，再令 \(y'=(1+x')/y=(x+1+y)/(xy)\)。每走一步，新的变量都是旧变量的有理式，分母似乎会越来越复杂；交替继续下去，出现的簇变量却只有五种。这种小规模的周期现象可以直接核对，但一般簇变异是否总能化成初始变量的 Laurent 多项式，需要更深的理由。另一方面，给 \(\mathfrak{sl}_2\) 的关系引入参数 \(q\)，又会改变表示中的系数与张量规则。本附录分别从可计算的模型看这些变化，而不把不同理论之间的类比当作证明。

\ProseParagraphBreak
本附录需要第14—15章的 \(\mathfrak{sl}_2\) 与包络代数、第21章的 Hopf 代数，以及第二卷第12.2节的有限长度模、合成列与 Jordan--Hölder 定理。簇代数小例只用有理函数的恒等变形，一般 Laurent 现象的证明概要另用通用系数和唯一分解。其余证明围绕所选模型展开；一般量子群、簇代数和高阶范畴化的系统理论可参阅相应专著。
\ProseParagraphBreak
量子 \(\mathfrak{sl}_2\) 与参数版本可参阅\cite[Section 5.6]{EtingofEtAl2015TensorCategories}；簇代数原始文献为\cite[Sections 2--3, 6]{FominZelevinsky2002ClusterFoundations}，所引节号依作者公开稿。

% !TeX root = ../../Book5.tex
% 纲要标题：### B.1 量子群与 Hopf 代数
% 纲要位置：工作记录/计划纲要.md，第 4955 行。
% 纲要细目（保留原定写作范围）：
%
% #### B.1.1 生成关系与余乘法
% #### B.1.2 权模与量子整数
% #### B.1.3 张量次序与非余交换性
%
\section{量子群与 Hopf 代数}
\label{b5:sec:B01}

本节以 \(U_q(\mathfrak{sl}_2)\) 介绍量子群，验证它的 Hopf 结构，并比较参数对表示的影响。量子半单代数的一般理论及其分类可留待进一步学习。

\subsection{生成关系与余乘法}
\label{b5:subsec:B01:01}

\begin{definition}\label{b5:def:B-quantum-sl2}
设 \(k\) 为特征零域，\(q\in k^\times\)、\(q^2\ne1\)。由生成元 \(E,F,K,K^{-1}\) 及关系
\[
KK^{-1}=K^{-1}K=1,\quad KEK^{-1}=q^2E,\quad KFK^{-1}=q^{-2}F,\quad
EF-FE=\frac{K-K^{-1}}{q-q^{-1}}
\]
定义的含幺结合 \(k\) 代数，记为 \(U_q(\mathfrak{sl}_2)\)，称为相应的量子包络代数。
\end{definition}
\symidx{\(U_q(\mathfrak{sl}_2)\)：采用给定关系的量子包络代数}
为解释 \(K\) 与经典 Cartan 元 \(H\) 的联系，可以在形式参数 \(\hbar\) 的意义下写 \(q=e^\hbar\)、\(K=e^{\hbar H}\)。形式幂级数给出
\[
\frac{K-K^{-1}}{q-q^{-1}}
=\frac{\sinh(\hbar H)}{\sinh\hbar}
=H+\frac{\hbar^2}{6}(H^3-H)+O(\hbar^4).
\]
这里分子、分母先约去共同因子 \(\hbar\)，余下分母有非零常数项，故在形式幂级数环中可逆。类似地，\(e^{\hbar H}Ee^{-\hbar H}=e^{2\hbar}E\) 的一次项恢复 \([H,E]=2E\)，对 \(F\) 得到 \([H,F]=-2F\)。这说明指数化的 \(K\) 如何记录经典关系；它是解释动机的形式模型，不能在原来的含分母公式中直接把数值参数 \(q\) 代成一。
\begin{proposition}\label{b5:prop:B-quantum-hopf}
设 \(k\) 为特征零域，\(q\in k^\times\)、\(q^2\ne1\)，\(U=U_q(\mathfrak{sl}_2)\) 为由 \(E,F,K,K^{-1}\) 生成的含幺结合 \(k\) 代数，关系为
\[
KK^{-1}=K^{-1}K=1,\quad KEK^{-1}=q^2E,\quad KFK^{-1}=q^{-2}F,\quad
EF-FE=\frac{K-K^{-1}}{q-q^{-1}}.
\]
则以下公式延伸为代数同态 \(\Delta:U\to U\otimes_kU\)、\(\varepsilon:U\to k\) 和反代数同态 \(S:U\to U\)，并给出 \(U\) 的 Hopf 结构：
\[
\begin{aligned}
\Delta(K)&=K\otimes K,&
\Delta(E)&=E\otimes K+1\otimes E,&
\Delta(F)&=F\otimes1+K^{-1}\otimes F,\\
\varepsilon(K)&=1,&\varepsilon(E)&=0,&\varepsilon(F)&=0,\\
S(K)&=K^{-1},&S(E)&=-EK^{-1},&S(F)&=-KF.
\end{aligned}
\]
\end{proposition}
\begin{proof}
先检查余乘法在自由代数上定义的同态保持关系。关于 \(K\) 的逆及共轭等式逐项由 \(KE=q^2EK\)、\(KF=q^{-2}FK\) 得到。交换子中两项交叉乘积相消，因为
\(EK^{-1}=q^2K^{-1}E\)、\(KF=q^{-2}FK\)。故
\[
\bigl[\Delta(E),\Delta(F)\bigr]
=[E,F]\otimes K+K^{-1}\otimes[E,F]
=\frac{K\otimes K-K^{-1}\otimes K^{-1}}{q-q^{-1}},
\]
正是右侧关系的余乘法。余单位也保持关系。对极按反同态延拓；例如
\(S(E)S(K)-q^2\cdot S(K)S(E)=0\)
对应 \(KE=q^2EK\) 的反序像；\(F\) 的共轭关系同理由
\(KF=q^{-2}FK\) 给出。剩下的交换子关系为
\[
 S(F)S(E)-S(E)S(F)
 =KFEK^{-1}-EF
 =FE-EF
 =\frac{K^{-1}-K}{q-q^{-1}},
\]
正是 \(S\bigl((K-K^{-1})/(q-q^{-1})\bigr)\)；这里 \(K\) 与 \(FE\) 交换。逆元关系也被保持，故三映射在商代数上良定义。
余结合与余单位只需在 \(E,F,K^{\pm1}\) 上核验；例如两种三重余乘法在 \(E\) 上都为
\(E\otimes K\otimes K+1\otimes E\otimes K+1\otimes1\otimes E\)。
对极的左、右等式在 \(E\) 上分别为
\(-EK^{-1}K+E=0\)、\(EK^{-1}-EK^{-1}=0\)，在 \(F\) 上分别为
\(-KF+KF=0\)、\(F-K^{-1}KF=0\)，在 \(K^{\pm1}\) 上都是逆元等式。
还须说明生成元上的核验能推广到乘积。用
\(\Delta(a)=\sum a_{(1)}\otimes a_{(2)}\) 的记号，若左对极等式已对 \(a,b\) 成立，则
\[
\begin{aligned}
\sum S\bigl(a_{(1)}b_{(1)}\bigr)a_{(2)}b_{(2)}
&=\sum S\bigl(b_{(1)}\bigr)S\bigl(a_{(1)}\bigr)a_{(2)}b_{(2)}\\
&=\varepsilon(a)\sum S\bigl(b_{(1)}\bigr)b_{(2)}
=\varepsilon(a)\varepsilon(b)1.
\end{aligned}
\]
右等式按相反次序先消去 \(b\) 的两项，得到同一结论。由乘积归纳及线性性，全部对极等式成立。
\end{proof}
这些公式及其版本可对照\cite[Definition 5.6.1; Theorem 5.6.2]{EtingofEtAl2015TensorCategories}。同一量子代数有不同但等价的余乘法惯例，计算张量表示时必须选定一套，不能混合来自两套惯例的公式。
\ProseParagraphBreak
余乘法也发生了变形，因为它规定同一个元素怎样作用于两个表示的张量积。在上面的形式模型中，
\[
\begin{aligned}
\Delta_\hbar(E)&=E\otimes1+1\otimes E+\hbar E\otimes H+O(\hbar^2),\\
\Delta_\hbar(F)&=F\otimes1+1\otimes F-\hbar H\otimes F+O(\hbar^2).
\end{aligned}
\]
常数项恢复经典包络代数的余乘法 \(\Delta_0(X)=X\otimes1+1\otimes X\)。记 \(\tau(a\otimes b)=b\otimes a\) 为普通张量翻转。若以 \(\Delta_\hbar^{\mathrm{op}}=\tau\Delta_\hbar\) 表示翻转两个张量因子后的余乘法，则
\[
\begin{aligned}
\Delta_\hbar(E)-\Delta_\hbar^{\mathrm{op}}(E)
 &=\hbar(E\otimes H-H\otimes E)+O(\hbar^2),\\
\Delta_\hbar(F)-\Delta_\hbar^{\mathrm{op}}(F)
 &=\hbar(F\otimes H-H\otimes F)+O(\hbar^2).
\end{aligned}
\]
因此经典极限虽然余交换，一次项已经区分两个张量因子的次序。量子关系与量子余乘法分别改变单个表示中的升降作用和两个表示之间的张量作用；下一小节的模型可以把这些区别实际算出。

\subsection{权模与量子整数}
\label{b5:subsec:B01:02}

对整数 \(r\)，定义量子整数
\[
[r]_q=\frac{q^r-q^{-r}}{q-q^{-1}}.
\]
\symidx{\([r]_q\)：量子整数}
固定整数 \(m\ge0\)。从一个满足 \(Ev_0=0\)、\(Kv_0=q^mv_0\) 的最高权向量出发，令 \(v_i=F^iv_0\)。用 \(Ev_0=0\) 和交换关系递推 \(EF^iv_0\)，得到 \(Ev_i=c_iv_{i-1}\)，其中系数满足
\[
c_{i+1}-c_i=[m-2i]_q,\qquad c_0=0.
\]
因此 \(c_i=\sum_{j=0}^{i-1}[m-2j]_q\)。展开两段几何级数，或逐项验证该和的递推，得到 \(c_i=[i]_q[m-i+1]_q\)。特别地 \(c_{m+1}=0\)，说明令 \(v_{m+1}=0\) 时，\(E\) 不会把这一被截去的向量带回保留的链段。为得到恰有 \(m+1\) 维的模型，下面直接在一组抽象基上定义作用并验证关系；任意给定最高权向量的幂链本身未必线性无关。这与经典升降算子的系数产生方式相同。
对 \(m\ge0\)，在基 \(v_0,\ldots,v_m\) 上规定
\[
Kv_i=q^{m-2i}v_i,\qquad Fv_i=v_{i+1},\qquad
Ev_i=[i]_q[m-i+1]_qv_{i-1},
\]
其中两端之外的向量为零。在这组模型中，给 \(v_i\) 标记整数权 \(m-2i\)，相应的 \(K\)-特征值为 \(q^{m-2i}\)。\(E\)、\(F\) 分别使权增加二和减少二，所以关于 \(K\) 的两个共轭关系直接成立；交换子关系等价于
\[
[i+1]_q[m-i]_q-[i]_q[m-i+1]_q=[m-2i]_q,
\]
确实，通分后的分子为
\[
\begin{aligned}
 &(q^{i+1}-q^{-i-1})(q^{m-i}-q^{-m+i})\\
 &\quad-(q^i-q^{-i})(q^{m-i+1}-q^{-m+i-1})\\
 &=(q-q^{-1})(q^{m-2i}-q^{-m+2i}).
\end{aligned}
\]
除以 \((q-q^{-1})^2\)，便得到右端的 \([m-2i]_q\)。所以这些矩阵确实构成一个 \(U_q(\mathfrak{sl}_2)\) 模。
\ProseParagraphBreak
若 \(q\) 不是单位根，则这些 \(K\)-特征值互异，所有内部升算子系数非零。记 \(\lambda_i=q^{m-2i}\)，插值多项式 \(p_i(T)=\prod_{j\ne i}(T-\lambda_j)/(\lambda_i-\lambda_j)\) 使 \(p_i(K)\) 成为投到 \(kv_i\) 的投影。非零不变子空间对这些多项式仍稳定，故含某个 \(v_i\)。又 \(E^iv_i=(\prod_{j=1}^i[j]_q[m-j+1]_q)v_0\ne0\)，所以含 \(v_0\)；再由 \(F^jv_0=v_j\) 得到整组基，证明该模不可约。上述结论针对这组最高权模型；一般量子群模的分类还须规定相应的表示类别。

\subsection{张量次序与非余交换性}
\label{b5:subsec:B01:03}

在 \(m=1\) 的二维模型中，\(Ev_1=v_0\)、\(Ev_0=0\)、\(Kv_0=qv_0\)。由固定的余乘法，
\[
E(v_1\otimes v_0)=q\,v_0\otimes v_0,\qquad
E(v_0\otimes v_1)=v_0\otimes v_0.
\]
当 \(q\ne1\) 时，普通翻转 \(\tau(v\otimes w)=w\otimes v\) 不与 \(E\) 的作用交换，因而不是 \(U_q(\mathfrak{sl}_2)\)-模同态。

\begin{proposition}\label{b5:prop:B-quantum-hecke-operator}
设 \(k\) 为特征零域，\(q\in k^\times\)、\(q^2\ne1\)。令 \(V\) 为 \(U_q(\mathfrak{sl}_2)\) 的二维模，满足
\(Kv_0=qv_0\)、\(Kv_1=q^{-1}v_1\)、\(Ev_1=v_0\)、\(Fv_0=v_1\)，其余升降作用为零；张量作用采用 \(\Delta(E)=E\otimes K+1\otimes E\)、\(\Delta(F)=F\otimes1+K^{-1}\otimes F\)。依次按基 \(v_0\otimes v_0,v_0\otimes v_1,v_1\otimes v_0,v_1\otimes v_1\)，定义
\[
\check R=\begin{pmatrix}
q&0&0&0\\
0&0&1&0\\
0&1&q-q^{-1}&0\\
0&0&0&q
\end{pmatrix}.
\]
则 \(\check R\) 是可逆的 \(U_q(\mathfrak{sl}_2)\)-模同态，且满足
\[
(\check R-qI)(\check R+q^{-1}I)=0,
\qquad
\check R_{12}\check R_{23}\check R_{12}
=\check R_{23}\check R_{12}\check R_{23}
\]
于 \(V^{\otimes3}\)，其中 \(\check R_{12}=\check R\otimes I\)、\(\check R_{23}=I\otimes\check R\)。
\end{proposition}
\symidx{\(\check R\)：二维量子模的 Hecke 翻转算子}
\begin{proof}
记 \(a=q-q^{-1}\)。在所列基上，\(E,F\) 的张量作用矩阵为
\[
\Delta(E)=\begin{pmatrix}0&1&q&0\\0&0&0&q^{-1}\\0&0&0&1\\0&0&0&0\end{pmatrix},
\qquad
\Delta(F)=\begin{pmatrix}0&0&0&0\\q^{-1}&0&0&0\\1&0&0&0\\0&1&q&0\end{pmatrix}.
\]
将它们与 \(\check R\) 相乘，两种次序分别相同；所用恒等式是 \(1+qa=q^2\) 与 \(q^{-1}+a=q\)。而 \(\Delta(K)=\operatorname{diag}(q^2,1,1,q^{-2})\)，\(\check R\) 只混合中间两个特征值同为一的方向，故也与 \(K^{\pm1}\) 交换。生成元的相容性给出模同态。中间二阶块 \(\begin{psmallmatrix}0&1\\1&a\end{psmallmatrix}\) 的平方为 \(I+a\begin{psmallmatrix}0&1\\1&a\end{psmallmatrix}\)，两端也有 \(q^2=1+aq\)，所以 \(\check R^2=I+a\check R\)。这既给出二次关系，也给出 \(\check R^{-1}=\check R-aI\)。

核验辫关系时，写 \(e_{ijk}=v_i\otimes v_j\otimes v_k\)。两种三重复合在 \(e_{000},e_{111}\) 上都乘以 \(q^3\)；在有序基
\((e_{001},e_{010},e_{100})\) 上，以及在 \((e_{011},e_{101},e_{110})\) 上，两者都给出矩阵
\[
\begin{pmatrix}0&0&q\\0&q&qa\\q&qa&q^2a\end{pmatrix}.
\]
例如 \(e_{100}\) 的像为 \(qe_{001}+qae_{010}+q^2ae_{100}\)，计算中用 \(a+qa^2=q^2a\)。这覆盖了八个基向量，故两种复合相等。
\end{proof}

把附录 A 的 Hecke 参数 \(v\) 取为这里的 \(q\)，这正是其二次关系与辫关系。在 \(V^{\otimes n}\) 的相邻两因子上放置 \(\check R\)，相距超过一的算子也交换，因此得到相应的 Hecke 代数作用。\(\check R\) 的矩阵本身在 \(q=1\) 可代入，并成为普通翻转；这不等于在量子代数原有的分式关系中直接代入一。在 \(k=\mathbb C,q=i\) 时，二次关系变成 \((\check R-iI)^2=0\)，而中间块使 \(\check R-iI\ne0\)。于是出现非平凡 Jordan 块，与\cref{b5:prop:hecke-rank-one-specialization}的双数代数相呼应。
\ProseParagraphBreak
这个算例给出了所选二维模及其张量幂上的算子。要为整个表示范畴构造辫结构，还须在不同对象之间给出自然的交换同构，并核验它们与张量结合的相容性。一般 Hopf 代数的模范畴具有幺半结构，但未必具有辫结构；若所选表示类别上有满足相应相容恒等式的 \(R\) 矩阵作用，可以将它与翻转复合，构造辫结构。拟三角 Hopf 代数的一般理论见\cite[Section 8.3, Definition 8.3.1, pp. 198--199]{EtingofEtAl2015TensorCategories}。

% !TeX root = ../../Book5.tex
% 纲要标题：### B.2 簇代数的变异与小秩例子
% 纲要位置：工作记录/计划纲要.md，第 4963 行。
% 纲要细目（保留原定写作范围）：
%
% #### B.2.1 种子与交换关系
% #### B.2.2 二阶五周期的完整计算
% #### B.2.3 Laurent 现象与计算范围
%
\section{簇代数的变异与小秩例子}
\label{b5:sec:B02}

簇代数通过一串有理变量替换生成交换环。每次替换只改变一个变量，新变量的分子是两个单项式的和。这些替换不断引入分母，但在某些小秩例子中会出现有限循环，变量仍能写成 Laurent 多项式。下面从一个二阶例子说明这一现象。

\subsection{种子与交换关系}
\label{b5:subsec:B02:01}

先取特征零域 \(k\)、有理函数域 \(k(x_1,\ldots,x_n)\)，以及反对称整数矩阵 \(B=(b_{ij})\)。向量 \(\mathbf x=(x_1,\ldots,x_n)\) 的分量代数独立，这一对数据称为一个无系数种子。沿位置 \(r\) 的变异保留 \(x_i\)、\(i\ne r\)，并用
\[
	x_rx'_r=\prod_{b_{ir}>0}x_i^{b_{ir}} +\prod_{b_{ir}<0}x_i^{-b_{ir}}
\]
定义 \(x'_r\)。空乘积取一。矩阵改为
\[
	b'_{ij}=
	\begin{cases}
		-b_{ij},&i=r\;\text{或}\;j=r,\\
		b_{ij}+[b_{ir}]_+[b_{rj}]_+-[-b_{ir}]_+[-b_{rj}]_+,&i,j\ne r,
	\end{cases}
	\quad [a]_+=\max(a,0).
\]
全部有限次变异得到的变量称为簇变量，它们在原有理函数域中生成的 \(k\) 子代数称为此种子的簇代数。 对反对称矩阵，约定 \(b_{ij}>0\) 表示从 \(i\) 到 \(j\) 有 \(b_{ij}\) 条箭头。例如箭图 \(1\to2\to3\) 的矩阵满足 \(b_{12}=b_{23}=1,b_{13}=0\)。在顶点 \(2\) 处变异时，
\[
	b'_{13}=0+[b_{12}]_+[b_{23}]_+=1.
\]
新增的 \(1\to3\) 正是沿两步路径 \(1\to2\to3\) 增加的箭头；同时原来的两条箭头反向，得到 \(2\to1\)、\(3\to2\)。一般情形的乘积项数出经过变异顶点的两步路径，带相反方向的修正项则对应消去二循环。因而任意顶点的变异还包含增添与消去箭头，只有源点、汇点情形才只需反向。新箭头还会改变下一次变量变异的交换单项式。
\symidx{\([a]_+\)：簇变异公式中的非负部分}

\begin{proposition}\label{b5:prop:B-three-vertex-mutations}
	设 \(k\) 为特征零域，\(x_1,x_2,x_3\) 在 \(k\) 上代数独立。对无系数种子及交换矩阵
	\[
		B=\begin{pmatrix}0&1&0\\-1&0&1\\0&-1&0\end{pmatrix},
	\]
	依次沿位置 \(2,1\) 变异，所得变量为
	\[
		x'_2=\frac{x_1+x_3}{x_2},\qquad x'_1=\frac{x'_2+x_3}{x_1},
	\]
	而第一次变异后的矩阵为
	\[
		B'=\begin{pmatrix}0&-1&1\\1&0&-1\\-1&1&0\end{pmatrix}.
	\]
\end{proposition}
\begin{proof}
	初始第二列的正、负条目分别在第一、第三行，故 \(x_2x'_2=x_1+x_3\)。矩阵第二行、第二列反号，\(b'_{13}=1\)、\(b'_{31}=-1\)，得到所列 \(B'\)。下一步第一列的正条目为 \(b'_{21}=1\)，负条目为 \(b'_{31}=-1\)，所以 \(x_1x'_1=x'_2+x_3\)。第二式中的 \(x_3\) 来自刚增加的箭头 \(1\to3\)；若只反向原有箭头，便会错误地把这一项写成空乘积一。
\end{proof}
\ProseParagraphBreak
两次在相同位置变异会回到原种子。因为第 \(r\) 列变号，两交换单项式互换，分子不变，故 \(x''_r=x_r\)。矩阵的两项修正也互换符号，相消后恢复原值；同时 \(b'_{ji}=-b'_{ij}\)，所以反对称性保持。这个可逆有理替换还说明新变量仍然代数独立并生成同一有理函数域。 这里采用无冻结变量、反对称矩阵的基本版本；带系数、冻结变量和反对称可对称化矩阵的其他版本还需补充相应定义。

\subsection{二阶五周期的完整计算}
\label{b5:subsec:B02:02}

取 \(B=\begin{psmallmatrix}0&1\\-1&0\end{psmallmatrix}\)，沿两个位置交替变异。令连续簇变量为 \(z_1=x,z_2=y\)，则
\(z_{i-1}z_{i+1}=1+z_i\)。 直接计算得
\[
	z_3=\frac{1+y}{x},\qquad z_4=\frac{1+x+y}{xy},\qquad z_5=\frac{1+x}{y},\qquad z_6=x,\quad z_7=y.
\]
例如 \(1+z_4=(x+1)(y+1)/(xy)\)，除以 \(z_3\) 得到 \(z_5\)；再继续两步就回到原来的两变量。任意变异串可消去相邻两次相同位置的变异；两个位置之间剩下的串必然交替，所以这个递推穷尽了所有可能。五个所列有理式两两不同，因此恰有五个可能簇变量，并有五个无序簇 \(\{z_i,z_{i+1}\}\)，下标模五。

\ProseParagraphBreak

无序簇的周期为五。若固定两个位置的编号，每次变异都把交换矩阵变为其相反矩阵；沿 \(1,2,1,2,1\) 变异五次，所得种子为
\[
	\Sigma_5=((y,x),-B).
\]
这正是把原种子的两个位置交换后得到的种子，因为交换行、列编号也把 \(B\) 变为 \(-B\)。继续沿 \(2,1,2,1,2\) 变异，十步后才恢复 \(\Sigma_{10}=((x,y),B)\)，所以有编号种子的周期为十。五个簇变量都为原变量的 Laurent 多项式，但一般不是普通多项式。

\ProseParagraphBreak

原始论文中的二阶递推见\cite[Example 2.5; Section 6]{FominZelevinsky2002ClusterFoundations}，本节选择了其中系数一、两个指数都为一的无系数 \(A_2\) 情形。
\ProseParagraphBreak
这个名称可以在已算出的分母中看到。\(A_2\) 的正根为 \(\alpha_1,\alpha_2,\alpha_1+\alpha_2\)，加上两个负简单根，得到几乎正根集合 \(\Phi_{\ge-1}=\Phi^+\cup\{-\alpha_1,-\alpha_2\}\)。将初始变量 \(z_1,z_2\) 分别标为 \(-\alpha_1,-\alpha_2\)，其余三个变量的约分后分母 \(x,xy,y\) 则分别给出正根 \(\alpha_1,\alpha_1+\alpha_2,\alpha_2\) 的简单根坐标 \((1,0),(1,1),(0,1)\)。这给出此算例的五个变量与五个几乎正根的对应，也说明根系记号为何出现在有理递推中。这里的五个簇与 \(A_2\) Weyl 群的六个 Weyl 室参数化的是不同对象，不能把两个数目混同；一般有限型的根系对应还需要额外的分类定理。

\subsection{Laurent 现象与计算范围}
\label{b5:subsec:B02:03}

\begin{theorem}[Laurent 现象]\label{b5:thm:B-laurent}
	设 \(k\) 为特征零域，\(n\ge1\) 为整数，\(x_1,\ldots,x_n\) 在 \(k\) 上代数独立，\(\mathcal F=k(x_1,\ldots,x_n)\)，\(B\in M_n(\mathbb Z)\) 为反对称交换矩阵。以 \(\Sigma=((x_1,\ldots,x_n),B)\) 为 \(\mathcal F\) 中没有系数、没有冻结变量的初始种子，按无系数种子变异产生的每个簇变量都属于 \(k[x_1^{\pm1},\ldots,x_n^{\pm1}]\)。
\end{theorem}
\symidx{\(\Sigma=(\mathbf x,B)\)：无系数簇种子及其交换矩阵}
\begin{proofsketch}
	证明不能只对路径长度直接使用交换式，因为除以一个已知 Laurent 多项式可能带来非单项式分母。解决这一点的办法，是把有限变异路径与每个路径顶点的邻边一起放进一棵有限树，并同时比较几条较短路径所得的表达。

	先把交换式的系数提升到通用系数的多项式环：固定初始交换多项式及下述相容关系中 \(C\) 的单项式支集，把这些多项式的系数暂作独立不定元，最后代入原交换多项式及相容因子的对应系数。后续交换多项式由相容关系递推产生，其系数不再另取独立参数。记这个系数环为 \(A\)，初始 Laurent 环为 \(L_0=A[x_1^{\pm1},\ldots,x_n^{\pm1}]\)。对连续标号 \(i,j,i\) 的三步变异，省略未改变的变量，写成
	\[
		z=\frac{P(y)}x,\qquad u=\frac{Q(z)}y,\qquad v=\frac{R(u)}z.
	\]
	这四个簇与通往目标簇 \(\Sigma_*\) 的路径可写成
	\[
		\begin{tikzcd}[column sep=large,row sep=large]
			(x,y)\arrow[r,"i"] &(z,y)\arrow[r,"j"] &(z,u)\arrow[r,"i"]\arrow[d,dashed,"\text{其余变异}"'] &(v,u)\\
			&&\Sigma_* &
		\end{tikzcd}
	\]
	省略的其余变量在前三条边上保持不变；右端 \((v,u)\) 是有限树中的一条支边，不必属于原来的变异路径。矩阵变异规定的指数变换给出交换多项式的相容关系 \(R(U)=C(U)P\bigl(Q(0)/U\bigr)\)，其中 \(C\) 对 \(U\) 是多项式，对未改变变量是 Laurent 多项式。将 \(v=R(Q(z)/y)/z\) 拆成其在 \(z=0\) 处的常数部分与差商；相容关系给出 \(R(Q(0)/y)=C(Q(0)/y)P(y)\)，再用 \(P(y)/z=x\)，便得
	\[
		v=\frac{R\bigl(Q(z)/y\bigr)-R\bigl(Q(0)/y\bigr)}z +xC\bigl(Q(0)/y\bigr).
	\]
	第一项的分子是关于 \(z\) 的多项式且常数项为零，所以能够除以 \(z\) 而不增加分母。这先证明三步后的变量仍在初始 Laurent 环中。

	在通用系数环中还可证明 \(\gcd(z,u)=\gcd(z,v)=1\)。由于 \(y\) 可逆，而 \(Q(z)-Q(0)\) 被 \(z\) 整除，有 \(\gcd(z,u)=\gcd(z,Q(z))=\gcd(z,Q(0))\)，这里最大公因数均只确定到单位。再由 \(z=P(y)/x\) 与 \(x\) 可逆，归结为 \(\gcd(P(y),Q(0))=1\)。例如把 \(P(y)\) 的一个系数参数记为 \(a\)，便可写成 \(P(y)=aM+P_0\)，其中 Laurent 单项式 \(M\) 可逆，且 \(P_0\) 不含 \(a\)。模掉 \(P(y)\) 就等于令 \(a=-M^{-1}P_0\)，所得商环是删去参数 \(a\) 后的多项式与 Laurent 多项式环，仍为整环。因此 \(P(y)\) 为素元，而系数与它独立的非零 \(Q(0)\) 在该商环中仍非零，不能被它整除。

	对后一式，差商模 \(z\) 等于在 \(z=0\) 处的导数，由链式法则得到
	\[
		v\equiv R'\bigl(Q(0)/y\bigr)Q'(0)/y +xC\bigl(Q(0)/y\bigr)\pmod z.
	\]
	右端关于可逆变量 \(x\) 为一次式，其系数不含 \(x\)。模掉 \(P(y)\) 后，\(Q(0)/y\) 仍非零；\(C\) 的系数参数独立，所以 \(C(Q(0)/y)\) 仍非零。上段已证明模掉 \(P(y)\) 后仍为整环，且 \(x\) 仍是独立的可逆变量。这个一次式的 \(x\) 系数非零，故它在该商环中非零。因为 \(z=P(y)/x\) 与 \(P(y)\) 生成同一理想，上面的同余式说明 \(v\) 在商环中也非零，即 \(P(y)\) 不整除 \(v\)，即 \(\gcd(z,v)=1\)。现在对有限树的路径长度归纳：从 \((z,y)\) 沿 \((z,u)\) 到 \(\Sigma_*\)，以及从 \((v,u)\) 沿 \((z,u)\) 到 \(\Sigma_*\)，都比从 \((x,y)\) 出发的路径短。归纳假设使目标变量 \(X\) 同时属于
	\[
		L_1=A[z^{\pm1},y^{\pm1},x_k^{\pm1}\ (k\ne i,j)],\qquad L_3=A[v^{\pm1},u^{\pm1},x_k^{\pm1}\ (k\ne i,j)].
	\]
	因为 \(y,x_k\) 在 \(L_0\) 中可逆，且 \(z,u,v\in L_0\)，这两个 Laurent 表达分别可写成
	\[
		X=\frac{f}{z^a}=\frac{g}{u^bv^c}, \qquad f,g\in L_0,\quad a,b,c\ge0.
	\]
	通分得到 \(fu^bv^c=gz^a\)。由 \(\gcd(z,u)=\gcd(z,v)=1\) 可知 \(z^a\) 与 \(u^bv^c\) 互素，而 \(L_0\) 是唯一分解环，所以 \(z^a\mid f\)，从而 \(X\in L_0\)。

	这里省略了通用系数的统一选取、相容关系的逐指数核验，以及有限树归纳的顶点记号；它们分别保证系数特殊化合法、三步消去可在每个位置实行、任意有限路径都被覆盖。通用系数与有限树归纳见\cite[Section 3, Theorem 3.2 and Lemma 3.3]{FominZelevinsky2002ClusterFoundations}；矩阵变异与交换关系的相容性见同文\cite[Section 4, Definition 4.2, Proposition 4.3 and Proposition 4.5]{FominZelevinsky2002ClusterFoundations}。归纳所得系数是通用参数的多项式，因此将参数代入对应系数后，分母仍只有初始变量的单项式，给出本节的无系数版本。当某交换矩阵列为零时，交换多项式为 \(1+1=2\)，合并后的常数系数参数应代入二。
\end{proofsketch}

三步相容式在前面的 \(A_2\) 例子中可以直接展开。取 \(P(Y)=1+Y\)、\(Q(Z)=1+Z\)、\(R(U)=1+U\)，则 \(Q(0)=1\)，并可取 \(C(U)=U\)，因为 \(U(1+U^{-1})=1+U\)。于是差商公式成为
\[
	v=\frac{\bigl(1+(1+z)/y\bigr)-\bigl(1+1/y\bigr)}{z}+\frac{x}{y} =\frac1y+\frac xy=\frac{1+x}{y}.
\]
这正是已经算出的 \(z_5\)。使用 \(Q(0)\) 是为了保留关于 \(z\) 的常数项；减去这一项后，分子便有因子 \(z\)，从而消去对新变量的除法。一般证明把这个局部机制放进同时比较多条路径的归纳中。
\ProseParagraphBreak
证明中的互素性必须在通用系数环中用完以后再作特殊化。例如 \(X-a\) 与 \(X-b\) 在 \(k[a,b,X]\) 中互素，代入 \(a=b=0\) 后却都成为 \(X\)。所以不能先代入参数，再沿用通用互素性。上面的顺序是先据此消去非单项式分母，得到系数为通用参数多项式的 Laurent 表达，然后特殊化这条已得到的恒等式；保留下来的是表达式及其初始单项式分母。
\ProseParagraphBreak
Laurent 性给出子代数的包含，并不保证初始变量的逆也属于簇代数。在 \(A_2\) 例子中，全部簇变量已穷尽，故 \(\mathcal A=k[z_1,\ldots,z_5]\subseteq k[x^{\pm1},y^{\pm1}]\)。在后一个环上代入 \(y=-1\)，五个生成元依次成为 \(x,-1,0,-1,-(1+x)\)，所以 \(\mathcal A\) 的像包含于 \(k[x]\)。若 \(x^{-1}\in\mathcal A\)，其像也应在 \(k[x]\) 中，但实际仍为 \(x^{-1}\)，矛盾。同理交换 \(x,y\) 可知 \(y^{-1}\notin\mathcal A\)。因此这里的包含严格；簇变量能用初始变量的逆写出，不等于这些逆元也由簇变量生成。
\ProseParagraphBreak
Laurent 性也不自动意味着只有有限多个簇变量。例如改成二阶递推 \(z_{i-1}z_{i+1}=1+z_i^2\)，则会产生无限多个不同的簇变量，尽管 Laurent 性仍然成立。一个直接核验是先取 \(z_1=z_2=1\)：递推得到 \(z_3=2,z_4=5\)，此后若 \(z_i>z_{i-1}>0\)，则 \(z_{i+1}=(1+z_i^2)/z_{i-1}>z_i\)。这些有理函数在该点的值互异，所以变量也互异。有限型分类需要控制全部变异，不是看到前几步没有新变量就能断定。类似地，正系数性、规范基、量子化和表示论实现分别是额外问题，不能都从 Laurent 现象本身推出。

% !TeX root = ../../Book5.tex
% 纲要标题：### B.3 表示论中的范畴化
% 纲要位置：工作记录/计划纲要.md，第 4971 行。
% 纲要细目（保留原定写作范围）：
%
% #### B.3.1 Grothendieck 群与两种加法
% #### B.3.2 从张量恒等式到函子
% #### B.3.3 分次与复形中的符号
%
\section{表示论中的范畴化}
\label{b5:sec:B03}

表示环把表示的直和与张量积记成类的加法与乘法；若还加入短正合列关系，不同同构类型也可能具有相同的类。回到表示、态射与函子以后，原来的整数恒等式可能来自更具体的直和分解或正合列。范畴化研究的正是这种提升；首先要说明取的是哪一种 Grothendieck 群，以及哪些信息在取群时被忘掉。

\subsection{Grothendieck 群与两种加法}
\label{b5:subsec:B03:01}

设 \(\mathcal A\) 为对象同构类构成集合的加性范畴。由对象符号 \([X]\) 生成，并加入 \([X\oplus Y]=[X]+[Y]\) 的交换群，称为分裂 Grothendieck 群。若 \(\mathcal A\) 还是交换范畴，即第四卷第4.2节定义的具有核、余核及相应像同构性质的加性范畴，则可以进一步对所有短正合列加入 \([M]=[L]+[N]\)，得到这里记为 \(G_0(\mathcal A)\) 的群。
\symidx{\(G_0(\mathcal A)\)：按短正合列关系定义的 Grothendieck 群}

\ProseParagraphBreak
两种关系不能混同。设 \(k\) 为特征 \(p\) 的域，考虑循环群 \(C_p\) 在 \(k\) 上的有限维表示。由 \(kC_p\cong k[t]/(t^p)\)，不可分解模为 \(M_d=k[t]/(t^d)\)、\(1\le d\le p\)。它们在分裂 Grothendieck 群中各给独立的基元，因为 Krull--Schmidt 分解唯一；而在 \(G_0\) 中，它们都满足 \([M_d]=d[M_1]\)，故只剩一个基元。这里 \(G_0\) 只记录简单合成因子的重数，分裂 Grothendieck 群仍可区分非分裂的不可分解对象。
\begin{proposition}\label{b5:prop:B-tensor-grothendieck}
	设 \(\mathcal A\) 为对象同构类构成集合的有限长度交换范畴，即每个对象都有有限合成列，配双加性且在每个变量上正合的张量积 \(\otimes:\mathcal A\times\mathcal A\to\mathcal A\)。记 \(G_0(\mathcal A)\) 为对对象同构类加入全部短正合列关系的 Grothendieck 群，\([X]\) 为对象 \(X\) 的类。则对 \(X,Y\in\mathcal A\)，乘法 \([X][Y]=[X\otimes Y]\) 在 \(G_0(\mathcal A)\) 上良定义；若有结合约束及单位对象，就得到相应的结合含幺环。其加法群以简单对象同构类为基。
\end{proposition}
\begin{proof}
	固定一个对象作张量保持短正合列，所以对象乘法尊重定义群时加入的关系，双加性便使乘法下降。结合与单位同构给出环的结合律及单位。第二卷第12.2节的合成列与 Jordan--Hölder 证明只用子对象、商对象及同构定理，因而适用于这里的交换范畴。每个对象沿合成列等于其简单因子之和；Jordan--Hölder 保证重数不依赖列，所以取各简单重数给出到自由交换群的逆映射。
\end{proof}
Hopf 代数的有限维模范畴中，张量在底层域上正合，便给出这种环。这里的“有限长度”不意味着只有有限多个简单类型，也不意味着整个范畴半单。基础定义及例子可对照\cite[Definition 1.5.8; Section 4.5]{EtingofEtAl2015TensorCategories}。

\subsection{从张量恒等式到函子}
\label{b5:subsec:B03:02}

取 \(\mathfrak{sl}_2\) 的有限维复表示范畴，全部对象完全可约。以 \(L_m\)、\(m\ge0\) 为简单对象，则
\[
	[L_1][L_m]=[L_{m+1}]+[L_{m-1}] \quad(m\ge1),\qquad [L_1][L_0]=[L_1].
\]
这是\cref{b5:thm:sl2-clebsch-gordan}在复数域上取一个张量因子为 \(L_1\) 所得的分解，在 Grothendieck 环中的像。令 \(t=[L_1]\)，便可递推地把每个 \([L_m]\) 写成 \(t\) 的一个首一 \(m\) 次整数多项式。不同简单类线性无关，所以这些首一递推多项式也保证
\[
	G_0\bigl(\operatorname{Rep}_{\mathrm{fd}}\mathfrak{sl}_2\bigr) \cong\mathbb Z[t].
\]
在范畴中，乘以 \(t\) 被提升为函子 \(F(V)=L_1\otimes V\)，其在简单对象上的作用就是上述直和分解。对象以外还有交织映射、自然变换以及张量结合约束；这些信息在只取环时看不见。因此给出一个环同构，并没有自动构造范畴等价；给出一串正整数，也没有自动指定相应的自然变换及其关系。
\ProseParagraphBreak
例如 \([L_2]=t^2-1\) 在对象层面来自 \(L_1\otimes L_1\simeq L_2\oplus L_0\)，在函子层面则给出
\[
	F^2\simeq (L_2\otimes-)\oplus\operatorname{Id}.
\]
具体地，翻转 \(\tau\) 与 \(\mathfrak{sl}_2\) 作用交换，投影 \((1+\tau)/2\)、\((1-\tau)/2\) 分别取出 \(\operatorname{Sym}^2L_1=L_2\) 与 \(\bigwedge^2L_1=L_0\)。将这两个投影张量任意对象的恒等映射，便得到与每个表示同态相容的自然投影。这是实际函子的直和分解；环中出现的减号只是在直和关系取 Grothendieck 类后移项得到的。
\ProseParagraphBreak
在非半单范畴中，相同的 Grothendieck 类甚至不能保证张量函子同构。取特征二的域 \(k\) 及 \(C_2\) 的有限维表示范畴，记平凡模为 \(M_1=k\)，正则模为 \(M_2=kC_2\)。由非分裂短正合列 \(0\to M_1\to M_2\to M_1\to0\)，有 \([M_2]=2[M_1]\)。两个正合函子
\[
	F_{\mathrm{reg}}(V)=M_2\otimes_k V,\qquad F_{\mathrm{triv}}(V)=V\oplus V
\]
在 \(G_0\) 上都诱导乘以二：前一函子用群的对角作用，且张量保持上述短正合列。但在 \(V=M_1\) 处，它们分别给出 \(M_2\) 和 \(M_1\oplus M_1\)。生成元在前者上的矩阵可取 \(\begin{psmallmatrix}1&1\\0&1\end{psmallmatrix}\)，在后者上则为恒等矩阵，故二者不同构，两个函子也不可能自然同构。被 \(G_0\) 忘掉的正是这个扩张是否分裂的信息。

\subsection{分次与复形中的符号}
\label{b5:subsec:B03:03}

设 \(k\) 为域。有限维 \(\mathbb Z\)-分次 \(k\)-向量空间的分裂 Grothendieck 环可用 Laurent 变量 \(v\) 记录次数平移：
\[
	[V]=\sum_{j\in\mathbb Z}(\dim V_j)v^j.
\]
张量积按次数相加，这个对应给出与 \(\mathbb Z[v,v^{-1}]\) 的环同构。每个分次片的维数都非负，而形式差允许负系数。
\ProseParagraphBreak
这也解释了范畴化为何常先换基。附录 A 的 Hecke 标准基满足 \(H_s^2=1+(v-v^{-1})H_s\)，其中的负系数不能直接作为对象直和分解的重数。改用 \(C'_s=H_s+v^{-1}\)，\cref{b5:prop:hecke-rank-one-specialization}给出
\[
	(C'_s)^2=(v+v^{-1})C'_s.
\]
在某个分次张量范畴中，若约定内部次数平移 \(\langle r\rangle\) 满足 \([X\langle r\rangle]=v^r[X]\)，那么表示 \(C'_s\) 的分次对象 \(B_s\) 若具有张量分解
\[
	B_s\otimes B_s\simeq B_s\langle1\rangle\oplus B_s\langle-1\rangle,
\]
取类后就得到这条规范基关系。两个 Laurent 单项式在这里对应两个次数平移后的对象。这个计算给出了寻找对象及自然分解的目标；正系数本身还没有构造出这样的范畴，也没有规定其中的态射和自然变换。

\ProseParagraphBreak
取有界的分次复形 \(C^\bullet\)，各 \(C^i\) 都有限维，且微分 \(d^i:C^i\to C^{i+1}\) 的内部次数为零，即 \(d^i(C^i_j)\subseteq C^{i+1}_j\)。于是圈 \(Z^i=\ker d^i\)、边界 \(B^i=\operatorname{im}d^{i-1}\) 和上同调 \(H^i(C^\bullet)\) 都继承内部的 \(\mathbb Z\)-分次，并有分次短正合列
\[
	\begin{aligned}
		0&\longrightarrow Z^i\longrightarrow C^i\longrightarrow B^{i+1}\longrightarrow0,\\
		0&\longrightarrow B^i\longrightarrow Z^i\longrightarrow H^i(C^\bullet)\longrightarrow0.
	\end{aligned}
\]
这两列逐内部次数分裂，故在上述分裂 Grothendieck 群中有 \([C^i]=[B^i]+[H^i(C^\bullet)]+[B^{i+1}]\)。乘以 \((-1)^i\) 后求和，边界项成对抵消，得到 \(\mathbb Z[v,v^{-1}]\) 中的 Euler 类恒等式
\[
	\sum_i(-1)^i[C^i]=\sum_i(-1)^i[H^i(C^\bullet)].
\]

\ProseParagraphBreak
第四卷第17—18章建立了映射锥、好三角和导出范畴；该卷附录 B.2 将复形的交错和称为 Euler 类，并讨论三角 Grothendieck 群。沿用这些约定，复形移位的负号还可以从三角关系直接读出。在有界分次复形的导出范畴中，对好三角 \(X\to Y\to Z\to X[1]\) 加入 \([Y]=[X]+[Z]\)，便得到相应的三角 Grothendieck 群。零映射 \(X\to0\) 的映射锥是 \(X[1]\)，给出好三角
\[
	X\longrightarrow0\longrightarrow X[1]\xrightarrow{1}X[1].
\]
因此 \(0=[X]+[X[1]]\)，即 \([X[1]]=-[X]\)。内部次数平移改变每一项的分次，而复形移位改变上同调次数；同时作 \(\langle1\rangle\) 与 \([1]\) 时，Euler 类相应乘以 \(-v\)。

\ProseParagraphBreak
于是有两种性质不同的符号：内部次数平移产生 \(v\)，复形移位产生负号。张量函子与内部平移分别在对象上实现乘法与分次；要进一步范畴化量子群或 Hecke 代数，还须构造生成函子之间的自然变换并验证相应关系。

% !TeX root = ../../Book5.tex
% 纲要标题：### B.4 参数特殊化与理论的适用范围
% 纲要位置：工作记录/计划纲要.md，第 4979 行。
% 纲要细目（保留原定写作范围）：
%
% #### B.4.1 单位根处的表示变化
% #### B.4.2 有理变异与点上的退化
%
\section{参数特殊化与理论的适用范围}
\label{b5:sec:B04}

量子参数和簇变量都能代入具体值，但代入可能使原来可逆的系数消失。要比较一般参数与特殊参数，须先确定在哪一个整系数或 Laurent 系数环上定义对象，再说明取商后哪些结论仍然成立。

\subsection{单位根处的表示变化}
\label{b5:subsec:B04:01}

在 \(U_q(\mathfrak{sl}_2)\) 的模型中，不可约性证明使用 \([i]_q\ne0\) 以及 \(K\) 的特征值互异。若 \(q\) 为奇数阶 \(\ell\ge3\) 的本原单位根，则 \([\ell]_q=0\)，不同整数权也可能给出相同的 \(K\)-特征值。升降系数的消失不仅使模型可约，还会产生不能拆成直和的扩张。

\begin{proposition}\label{b5:prop:B-root-of-unity-filtration}
	设 \(k\) 为特征零域，\(q\in k\) 为奇数阶 \(\ell\ge3\) 的本原单位根。令 \(V_\ell\) 为 \(U_q(\mathfrak{sl}_2)\) 在基 \(v_0,\ldots,v_\ell\) 上的模，作用为
	\[
		Kv_i=q^{\ell-2i}v_i,\qquad Fv_i=v_{i+1},\qquad Ev_i=[i]_q[\ell-i+1]_qv_{i-1},
	\]
	两端之外的向量取零。记 \(L_r\) 为同一公式中 \(m=r\) 的 \(r+1\) 维模型。则
	\[
		0\subset T=kv_\ell\subset W=\operatorname{span}_k\{v_1,\ldots,v_\ell\}\subset V_\ell
	\]
	是合成列，其中 \(T\cong L_0\)、\(W/T\cong L_{\ell-2}\)、\(V_\ell/W\cong L_0\)。短正合列 \(0\to T\to W\to L_{\ell-2}\to0\) 不分裂。
\end{proposition}
\begin{proof}
	因为 \(Ev_1=[1]_q[\ell]_qv_0=0\)，其余 \(E,F,K^{\pm1}\) 的作用也保持 \(W\)，所以 \(W\) 是子模。又 \(Ev_\ell=[\ell]_q[1]_qv_{\ell-1}=0\)、\(Fv_\ell=0\)、\(Kv_\ell=v_\ell\)，故 \(T\) 为平凡子模；商 \(V_\ell/W\) 由 \(v_0\) 的类生成，\(E,F\) 都作用为零而 \(K\) 作用为一，也是平凡模。

	在 \(W/T\) 中令 \(w_j=v_{j+1}+T\)，\(0\le j\le\ell-2\)。由 \([\ell-a]_q=-[a]_q\)，有
	\[
		\begin{aligned}
			Kw_j&=q^{\ell-2-2j}w_j,& Fw_j&=w_{j+1},\\
			Ew_j&=[j+1]_q[\ell-j]_q w_{j-1} =[j]_q[\ell-1-j]_q w_{j-1}.
		\end{aligned}
	\]
	边端仍取零，恰为 \(L_{\ell-2}\) 的公式。因 \(\ell\) 为奇数，\([a]_q=0\) 等价于 \(\ell\mid a\)，故这里所有内部升算子系数非零；\(K\) 的这些特征值也互异。沿用\subsecref{b5:subsec:B01:02}的插值投影，对互异的 \(\lambda_j=q^{\ell-2-2j}\) 取 \(p_j(T)=\prod_{i\ne j}(T-\lambda_i)/(\lambda_j-\lambda_i)\)。\(p_j(K)\) 使非零不变子空间包含某个 \(w_j\)；各内部升算子系数非零，故 \(E^jw_j\) 为 \(w_0\) 的非零倍，再由 \(F^iw_0=w_i\) 得到全部基，证明 \(W/T\) 简单。因此三个商都是简单模。

	若 \(W=T\oplus C\) 为子模直和，\(C\to W/T\) 为同构，必有 \(C\) 中 \(K\)-特征值 \(q^{\ell-2}\) 的非零向量。\(W\) 上的 \(\ell\) 个 \(K\)-特征值互异，该特征空间就是 \(kv_1\)，所以 \(v_1\in C\)。但 \(F^{\ell-1}v_1=v_\ell\ne0\)，使 \(C\cap T\ne0\)，矛盾。故该扩张不分裂。
\end{proof}

这条合成列记录了两端的平凡因子和中间的简单因子，非分裂性则说明它们的重数还不足以确定整个模。这个最高权模型展示了单位根处的非半单现象，与 \subsecref{b5:subsec:B03:01}中特征 \(p\) 的群表示相呼应。
\ProseParagraphBreak
把 \(q\) 取为一时，不能直接代入 \((K-K^{-1})/(q-q^{-1})\)，因为分母变为零。须先选择适当的整形式或引入额外生成元，再讨论特殊化。这里整形式指在 \(\mathbb Z[q,q^{-1}]\)（或适当的系数子环）上定义、扩张到 \(\mathbb Q(q)\) 后恢复原代数的模型；它给出可按系数环同态 \(q\mapsto q_0\) 取张量的代数。若模型含需要取逆的系数，还须核对其像非零，不能把含 \((q-q^{-1})^{-1}\) 的模型直接代到 \(q_0=1\)。不同整形式及其单位根特殊化见\cite[Section 5.6, pp. 101--102]{EtingofEtAl2015TensorCategories}；\(q\to1\) 的形式极限另见该书 Remark 5.6.5。形式极限本身不足以确定特殊化后的代数。

\subsection{有理变异与点上的退化}
\label{b5:subsec:B04:02}

簇变异是一项可逆有理坐标变换，不是每个仿射点上都定义的正则同构。前面的 \(A_2\) 例子中，若取 \(x=y=-1\)，五个 Laurent 表达均有意义，得到
\[
	(z_1,z_2,z_3,z_4,z_5)=(-1,-1,0,-1,0).
\]
交换恒等式仍成立，但若从含 \(z_3=0\) 的簇出发，用除以 \(z_3\) 的公式恢复邻近变量，就不能直接在这个点上操作。一个在有理函数域中合法的逆变换，可能在特定点失去意义。
\ProseParagraphBreak
还可以看到，初始变量的取值未必确定簇代数的点。令 \(\mathcal A=k[z_1,\ldots,z_5]\) 为上述 \(A_2\) 簇代数，我们希望趋近 \((x,y)=(0,-1)\) 时，\(z_3=(1+y)/x\) 仍能取任意指定值 \(c\)。令 \(x=t\)、\(1+y=ct\) 正好使该比值恒为 \(c\)，故对任意 \(c\in k\) 考虑这族代入 \(x=t\)、\(y=-1+ct\)。初始 Laurent 环要求 \(x,y\) 及其逆都有像；因为 \(x\mapsto t\)、\(y\mapsto-1+ct=-(1-ct)\)，须先把 \(t\) 与 \(1-ct\) 都取逆。这便定义从初始 Laurent 环到 \(k[t^{\pm1},(1-ct)^{-1}]\) 的同态；五个生成元的像为
\[
	z_1=t,\qquad z_2=-1+ct,\qquad z_3=c,\qquad z_4=-\frac{1+c}{1-ct},\qquad z_5=-\frac{1+t}{1-ct}.
\]
初始 Laurent 环含有 \(x^{-1}\)，不能直接在它的像中令 \(t=0\)。但这些生成元的像全部落在子环 \(k[t,(1-ct)^{-1}]\) 中，所以限制到 \(\mathcal A\) 后可以继续取 \(t=0\)，得到一个真正的 \(k\) 代数同态 \(\mathcal A\to k\)，其五个取值为
\[
	(0,-1,c,-1-c,-1).
\]
不同 \(c\) 给出不同的点，但它们的初始坐标都为 \((x,y)=(0,-1)\)。在这个初始坐标处，\((1+y)/x\) 无法直接求值；作为 \(\mathcal A\) 的元素，\(z_3\) 却在每个上述点上都有值，而且该值由点中额外的坐标信息确定。因而应区分初始簇给出的非零坐标图与整个 \(\operatorname{Spec}\mathcal A\)：后者还包括不在这张图中的点。
\ProseParagraphBreak
计算时应区分这两种情形：变量作为代数独立有理函数时，变异可逆且满足已证恒等式；变量代入一个点时，需要逐项检查所用分母。类似的区分也适用于量子参数的特殊化和范畴化中的基变换。进一步学习时，还需分别讨论规范基、正性、量子簇结构和高阶范畴化的定义与定理。

% !TeX root = ../../Book5.tex
% 纲要标题：### B.5 辫张量结构与融合范畴
% 纲要位置：工作记录/计划纲要.md（2026-10-20 新增骨架）。
% 写作状态：2026-10-20 新增骨架；正文、证明、算例与习题均待撰写。
% 写作范围：复用 B.1 的二维 R 算子，刚性回引第四卷第5章；Rep(G) 与 Vec_G 固定有限群及特征零条件，Tannaka 回引附录 E.4。
% #### B.5.1 拟三角结构与自然辫
% #### B.5.2 刚性与融合范畴
% #### B.5.3 有限群表示与分次向量空间
% #### B.5.4 中性 Tannaka 重构与微分模

\section{辫张量结构与融合范畴}
\label{b5:sec:B05}

待撰写。

\subsection{拟三角结构与自然辫}
\label{b5:subsec:B05:01}

待撰写。

\subsection{刚性与融合范畴}
\label{b5:subsec:B05:02}

待撰写。

\subsection{有限群表示与分次向量空间}
\label{b5:subsec:B05:03}

待撰写。

\subsection{中性 Tannaka 重构与微分模}
\label{b5:subsec:B05:04}

待撰写。

% 新增结构与纲要同步；以下条目均为待撰写计划。
% * 状态：待撰写；当前仅有结构占位。

% !TeX root = ../../Book5.tex
% 纲要标题：### B.6 Hall 代数与箭图
% 纲要位置：工作记录/计划纲要.md（2026-10-20 新增骨架）。
% 写作状态：正文、证明、算例及配套习题与解答已完成初稿；新增内容尚未正式审核。
% 写作范围：在有限域上展开计数及 A1、A2 算例；准确区分 Hall 代数、合成子代数与量子群正部，不展开典范基及晶体基。
% #### B.6.1 Hall 数与结合乘法
% #### B.6.2 Euler 型与扭曲乘法
% #### B.6.3 单顶点与双顶点算例
% #### B.6.4 量子群正部与 Ringel 定理

\section{Hall 代数与箭图}
\label{b5:sec:B06}

给定两个模，短正合列说明它们可以怎样组成一个较大的模。若底域有限，还能数出某个中项中有多少个指定类型的子模，并把这个数作为乘法系数。三个模依次组成一个模时，两种加括号方式数的是同一批旗，这正是 Hall 乘法的结合律。对箭图再加入 Euler 型给出的系数，简单模之间的乘积便出现量子 Serre 关系。

\ProseParagraphBreak
本节用 \(q\) 表示有限域的元素个数，固定正实数 \(v=\sqrt q\)，以 \(K=\mathbb Q(v)\) 为计数代数的系数域。因此表示本身定义在 \(\mathbb F_q\) 上，Hall 代数却是特征零域 \(K\) 上的代数。\subsecref{b5:subsec:B01:02}中的量子参数在这里改记为 \(v\)；它是一个固定的数，区别于域的大小 \(q=v^2\)。

\subsection{Hall 数与结合乘法}
\label{b5:subsec:B06:01}

先不要求代数来自箭图。设 \(A\) 是有限域上的有限维代数，以下只考虑有限维左 \(A\)-模。子模的个数有限，且每个固定维数上只有有限多个模的同构类：取 \(A\) 的一组有限基，作用由有限域上的有限组矩阵给定，可能的矩阵组本身就只有有限多个。

\begin{definition}\label{b5:def:B-hall-number}
	设 \(k=\mathbb F_q\)，\(A\) 为有限维含幺 \(k\) 代数，\(M,N,L\) 为有限维左 \(A\)-模。定义\term[Hall-shu]{Hall 数}[Hall number]
	\[
		F_{M,N}^{L}
		=\#\{\,U\subseteq L\mid U\cong N,\ L/U\cong M\,\}.
	\]
	它数的是实际的子模 \(U\)，不另选择 \(N\to U\) 或 \(L/U\to M\) 的同构。令 \(v=\sqrt q>0\)、\(K=\mathbb Q(v)\)，以模的同构类为基，定义
	\[
		\mathcal H_0(A)=\bigoplus_{[M]}K[M],\qquad
		[M]\circ[N]=\sum_{[L]}F_{M,N}^{L}[L].
	\]
	这里的和按 \(K\) 双线性延拓，所得乘法称为未扭曲的 Hall 乘法。
\end{definition}
\symidx{\(F_{M,N}^{L}\)：子模为 \(N\)、商模为 \(M\) 的 Hall 数}
\symidx{\(\mathcal H_0(A)\)：有限域上有限维代数的未扭曲 Hall 代数}
同构 \(L\to L'\) 给出相应子模集合的双射，所以这些数只依赖三个模的同构类。若 \(F_{M,N}^{L}\ne0\)，则 \(\dim_kL=\dim_kM+\dim_kN\)；固定这一维数只有有限多个同构类，因此乘积确实仍是有限线性组合。

\begin{proposition}\label{b5:prop:B-hall-associativity}
	设 \(k=\mathbb F_q\)、\(A\) 为有限维含幺 \(k\) 代数，\(K=\mathbb Q(\sqrt q)\)。\cref{b5:def:B-hall-number}中的乘法使 \(\mathcal H_0(A)\) 成为以 \([0]\) 为单位元的结合 \(K\) 代数。对任意有限维左 \(A\)-模 \(M,N,R,L\)，有
	\begin{equation}\label{b5:eq:B-hall-flag-identity}
		\sum_{[X]}F_{M,N}^{X}F_{X,R}^{L}
		=\sum_{[Y]}F_{N,R}^{Y}F_{M,Y}^{L}.
	\end{equation}
\end{proposition}
\begin{proof}
	固定 \(L\)，数其中满足
	\[
		0\subseteq U\subseteq V\subseteq L,
		\qquad U\cong R,\quad V/U\cong N,\quad L/V\cong M
	\]
	的旗。先选 \(U\)，并按 \(L/U\) 的同构类 \([X]\) 分组；每个这样的 \(U\) 有 \(F_{M,N}^{X}\) 个可选的商子模 \(V/U\)。子模对应定理将它们一一提升为包含 \(U\) 的子模 \(V\)，得到左边。先选 \(V\)，按其同构类 \([Y]\) 分组，再在 \(V\) 内选 \(U\)，则得到右边。计数中不选择商模到代表元的同构，因此没有额外的自同构群因子。

	等式两边分别是 \(([M]\circ[N])\circ[R]\) 与 \([M]\circ([N]\circ[R])\) 中 \([L]\) 的系数，故乘法结合。最后，商模为零的子模只能是 \(L\) 本身，子模为零时商只能是 \(L\)，所以 \([0]\circ[M]=[M]=[M]\circ[0]\)。
\end{proof}
同一论证给出任意多个因子的含义：\([M_1]\circ\cdots\circ[M_r]\) 中 \([L]\) 的系数，等于旗
\[
	0=L_r\subseteq L_{r-1}\subseteq\cdots\subseteq L_0=L,
	\qquad L_{i-1}/L_i\cong M_i\quad(1\le i\le r)
\]
的数目。因此右边的因子处在旗的下层。文献也常先数固定两端和中项的短正合列，再除以 \(|\operatorname{Aut}(M)|\cdot|\operatorname{Aut}(N)|\)，得到同一 Hall 数；这同直接数 \(\operatorname{Ext}^1_A(M,N)\) 的元素有区别。上述约定与旗计数证明可对照\cite[作者公开稿, Section 1.3, Proposition 1.1, Lemma 1.2]{Schiffmann2012HallAlgebras}。

\subsection{Euler 型与扭曲乘法}
\label{b5:subsec:B06:02}

现在取 \(A=kQ\)，其中 \(k=\mathbb F_q\)，\(Q\) 为有限无有向环箭图。路径代数有限维，故上一小节适用。令 \(\mathbf d(M)\) 为维数向量；短正合列逐顶点正合，因而乘积按维数向量相加分次。

\ProseParagraphBreak
\cref{b5:prop:quiver-euler-hom-ext}对任意域上的有限无环箭图给出
\begin{equation}\label{b5:eq:B-hall-euler}
	\begin{aligned}
	\langle\mathbf d,\mathbf e\rangle_Q
	&=\sum_{i\in Q_0}d_ie_i-\sum_{a:i\to j}d_ie_j,\\
	\langle\mathbf d(M),\mathbf d(N)\rangle_Q
	&=\dim_k\operatorname{Hom}_{kQ}(M,N)
	-\dim_k\operatorname{Ext}^1_{kQ}(M,N).
	\end{aligned}
\end{equation}
这个整数由两条维数向量决定，且对每个变量都是加性的。用它修正 Hall 乘法，不会破坏结合律。

\begin{definition}\label{b5:def:B-twisted-hall}
	设 \(Q\) 为有限无有向环箭图，\(k=\mathbb F_q\)、\(v=\sqrt q>0\)、\(K=\mathbb Q(v)\)。在 \(\mathcal H_0(kQ)\) 的同一 \(K\) 向量空间上，定义
	\begin{equation}\label{b5:eq:B-twisted-hall-product}
		[M]*[N]
		=v^{\langle\mathbf d(M),\mathbf d(N)\rangle_Q}
		\sum_{[L]}F_{M,N}^{L}[L].
	\end{equation}
	所得代数记为 \(\mathcal H_v(Q)\)，称为 \(Q\) 的\term[niuqu-Hall-daishu]{扭曲 Hall 代数}[twisted Hall algebra]。
\end{definition}
\symidx{\(\mathcal H_v(Q),\ *\)：箭图的扭曲 Hall 代数及其乘法}

\begin{proposition}\label{b5:prop:B-hall-twist-associativity}
	设 \(Q\) 为有限无有向环箭图，\(k=\mathbb F_q\)、\(v=\sqrt q>0\)。乘法\eqref{b5:eq:B-twisted-hall-product}使 \(\mathcal H_v(Q)\) 成为以 \([0]\) 为单位元的 \(\mathbb N^{Q_0}\) 分次结合 \(\mathbb Q(v)\) 代数。
\end{proposition}
\begin{proof}
	记三个因子的维数向量为 \(\mathbf d,\mathbf e,\mathbf f\)。\(([M]*[N])*[R]\) 的每个非零项有相同的附加系数
	\[
		v^{\langle\mathbf d,\mathbf e\rangle_Q+
		\langle\mathbf d+\mathbf e,\mathbf f\rangle_Q}.
	\]
	另一种括号给出的指数为 \(\langle\mathbf e,\mathbf f\rangle_Q+\langle\mathbf d,\mathbf e+\mathbf f\rangle_Q\)。双线性使两者都等于
	\[
		\langle\mathbf d,\mathbf e\rangle_Q+
		\langle\mathbf d,\mathbf f\rangle_Q+
		\langle\mathbf e,\mathbf f\rangle_Q.
	\]
	除去这个共同系数后，结合律就是\eqref{b5:eq:B-hall-flag-identity}。零向量与每个向量配对都为零，故单位元不变。每个中项的维数向量是两端之和，这也证明分次相容性。
\end{proof}
令 \(\alpha_i\) 为顶点 \(i\) 的标准单位向量，\(S_i\) 为相应简单模，则
\[
	\langle\alpha_i,\alpha_j\rangle_Q
	=\delta_{ij}-\#\{\,a:i\to j\,\}.
\]
尤其在本书的左模约定下，箭 \(1\to2\) 使 \(\langle\alpha_1,\alpha_2\rangle_Q=-1\)，反向配对为零。这决定下面两个乘积中的哪一个应带 \(v^{-1}\)。Euler 扭曲采用的是有序的双线性型；若只用其对称部分，将得不到同一套乘法系数。Euler 型及这一扭曲的定义见\cite[作者公开稿, Sections 1.2--1.3]{Schiffmann2012HallAlgebras}。

\subsection{单顶点与双顶点算例}
\label{b5:subsec:B06:03}

先取一个顶点且没有箭。一个表示只是有限维 \(k=\mathbb F_q\) 向量空间，令 \(u_r=[k^r]\)、\(u_0=[0]\)。在 \(k^d\) 中选一个 \(r\) 维子空间，先选它的一组有序基，再除以同一子空间的有序基数，得到
\[
	\#\operatorname{Gr}(r,k^d)
	=\frac{\prod_{j=0}^{r-1}(q^d-q^j)}{\prod_{j=0}^{r-1}(q^r-q^j)}
	=\prod_{j=1}^{r}\frac{q^{d-r+j}-1}{q^j-1}
	\coloneqq\begin{bmatrix}d\\r\end{bmatrix}_q.
\]
这个整数称为\term[Gauss-erxiangshi-xishu]{高斯二项式系数}[Gaussian binomial coefficient]；空积取一。所有短正合列都分裂，但子空间通常不止一个，因此即使中项只有一种，也会出现非平凡的计数系数。

\begin{proposition}\label{b5:prop:B-hall-one-vertex}
	设 \(Q\) 只有一个顶点且没有箭，\(k=\mathbb F_q\)、\(v=\sqrt q>0\)、\(u_r=[k^r]\)。则对 \(m,n\ge0\)，有
	\[
		u_m\circ u_n=\begin{bmatrix}m+n\\n\end{bmatrix}_q u_{m+n},
		\qquad
		u_m*u_n=v^{mn}\begin{bmatrix}m+n\\n\end{bmatrix}_q u_{m+n}.
	\]
	以 \([j]_v=(v^j-v^{-j})/(v-v^{-1})\)、\([r]_v!=\prod_{j=1}^{r}[j]_v\) 记量子整数及量子阶乘，则
	\begin{equation}\label{b5:eq:B-hall-divided-power}
		u_1^{*r}=v^{r(r-1)}[r]_v!\,u_r,
		\qquad \frac{u_1^{*r}}{[r]_v!}=v^{r(r-1)}u_r.
	\end{equation}
	所以 \(\mathcal H_v(Q)\cong K[T]\)，其中 \(K=\mathbb Q(v)\)，同构把 \(T\) 送到 \(u_1\)。
\end{proposition}
\begin{proof}
	Hall 数恰为上述子空间数；Euler 型在两个维数上为 \(mn\)，给出两种乘积。逐次与 \(u_1\) 相乘，得到
	\[
		u_1^{*r}=v^{r(r-1)/2}
		\prod_{j=1}^{r}\frac{q^j-1}{q-1}\,u_r.
	\]
	由 \(q=v^2\)，每个因子满足 \((q^j-1)/(q-1)=v^{j-1}[j]_v\)，再把这些幂相乘即得\eqref{b5:eq:B-hall-divided-power}。所有系数均非零，故 \(u_1\) 生成整个代数；不同次幂是不同基元素 \(u_r\) 的非零倍数，因而线性无关。这同时证明满射 \(K[T]\to\mathcal H_v(Q)\) 单射。
\end{proof}
这里 \([r]_v\) 是对称的量子整数，而 \((q^r-1)/(q-1)=1+q+\cdots+q^{r-1}\) 是有限域子空间计数中的整数；两者相差 \(v^{r-1}\)。这一个因子解释了除幂同 \([k^r]\) 之间的归一化。上述单顶点计算可对照\cite[作者公开稿, Section 3.3, Example 3.14]{Schiffmann2012HallAlgebras}。

\ProseParagraphBreak
现在取 \(Q:1\to2\)。记
\[
	\begin{gathered}
		S_1=(k\to0),\qquad S_2=(0\to k),\qquad P=(k\xrightarrow{1}k),\\
		D=S_1\oplus S_2,\qquad x=[S_1],\quad y=[S_2].
	\end{gathered}
\]
维数向量为 \((1,1)\) 的表示，按唯一箭的秩只分为 \(D\) 与 \(P\) 两类。在每一类中，顶点 \(2\) 的整个空间都给出唯一的子表示 \(S_2\)，商为 \(S_1\)。反过来，子表示若为 \(S_1\)，箭必须为零，所以只有 \(D\) 可以出现。于是
\begin{equation}\label{b5:eq:B-hall-a2-product}
	\begin{aligned}
	 x\circ y&=[D]+[P],& y\circ x&=[D],\\
	 x*y&=v^{-1}([D]+[P]),& y*x&=[D].
	\end{aligned}
\end{equation}
这同 \(\operatorname{Ext}^1(S_1,S_2)\cong k\)、\(\operatorname{Ext}^1(S_2,S_1)=0\) 一致，但非零扩张类的 \(q-1\) 个参数并没有成为 \([P]\) 的 Hall 系数：这里实际计数的子表示只有一个。这个算例也给出
\[
	[P]=v(x*y)-y*x.
\]
因此两个简单模的加权交换子已经恢复第三个不可分解表示。

\begin{proposition}\label{b5:prop:B-hall-a2-serre}
	设 \(Q:1\to2\)，\(k=\mathbb F_q\)、\(v=\sqrt q>0\)。在 \(\mathcal H_v(Q)\) 中，令 \(x=[S_1]\)、\(y=[S_2]\)，并在下式中省略扭曲乘号 \(*\)。则
	\begin{equation}\label{b5:eq:B-hall-a2-serre}
		\begin{aligned}
		x^2y-(v+v^{-1})xyx+yx^2&=0,\\
		y^2x-(v+v^{-1})yxy+xy^2&=0.
		\end{aligned}
	\end{equation}
\end{proposition}
\begin{proof}
	记 \(L_0=S_1^{\oplus2}\oplus S_2\)、\(L_1=S_1\oplus P\)。它们是维数向量 \((2,1)\) 的全部表示类型，分别对应零箭与秩一箭。单顶点计算给出 \(x^2=v(q+1)[S_1^{\oplus2}]\)。先取唯一的终点空间作为 \(S_2\) 子模，再用 Euler 指数 \(\langle2\alpha_1,\alpha_2\rangle_Q=-2\)，得到
	\[
		x^2y=v^{-1}(q+1)([L_0]+[L_1]).
	\]
	在 \(yx^2\) 中，所要求的子表示在顶点 \(1\) 占满 \(k^2\)、在顶点 \(2\) 为零，迫使箭为零；这时子表示唯一，Euler 指数为零，故
	\[
		yx^2=v(q+1)[L_0].
	\]
	最后由 \(yx=[D]\)，计算 \(x[D]\)。在 \(L_0\) 中，子表示 \(D\) 的起点可取 \(k^2\) 的任意直线，共 \(q+1\) 条；在 \(L_1\) 中，这条直线必须是秩一箭的核，只有一条。此时 Euler 指数 \(\langle\alpha_1,\alpha_1+\alpha_2\rangle_Q=0\)，所以
	\[
		xyx=(q+1)[L_0]+[L_1].
	\]
	分别比较 \([L_0]\)、\([L_1]\) 的系数，并用 \(v^{-1}(q+1)=v+v^{-1}\)，得到第一条关系。

	对表示作有限维向量空间对偶，再交换顶点 \(1,2\)，仍得到 \(Q\) 的表示。它反转短正合列，把子模与商模交换；同时反转 Euler 配对的次序。因此它在 Hall 代数上给出反同构，并交换 \(x,y\)。将这个反同构作用于第一条关系，所得为 \(xy^2-(v+v^{-1})yxy+y^2x=0\)，即第二条关系。
\end{proof}
若两个顶点之间没有边，则只有分裂扩张，两个方向的子表示都唯一，交叉 Euler 指数也都为零，故两个简单模的类交换。这一交换关系与\eqref{b5:eq:B-hall-a2-serre}正好涵盖单边 Dynkin 图中不同顶点的两种关系。

\subsection{量子群正部与 Ringel 定理}
\label{b5:subsec:B06:04}

设 \(Q\) 的无向基础图是连通的 \(A_n,D_n,E_6,E_7,E_8\) 型 Dynkin 图，\(\mathfrak g\) 为相应的复单 Lie 代数。这些图都只有单边，即两个相邻顶点之间恰连一条边。以顶点标记简单根 \(\alpha_i\)，Cartan 矩阵为
\[
	c_{ii}=2,\qquad
	c_{ij}=\begin{cases}-1,&i,j\text{ 相邻},\\0,&i,j\text{ 不相邻}\end{cases}
	\quad(i\ne j).
\]
它也是 Euler 型对称化后在标准基上的矩阵。量子正部的关系只用这个无向图；用于实现这些关系的 Hall 计数则保留了具体箭向。

\begin{definition}\label{b5:def:B-quantum-positive-part}
	设 \(q\) 为素数幂，\(v=\sqrt q>0\)、\(K=\mathbb Q(v)\)，\(\mathfrak g\) 为一个连通的单边 Dynkin 图所对应的复单 Lie 代数，以 \(I\) 为该图的顶点集。由生成元 \(E_i\)、\(i\in I\)，及以下关系定义的含幺结合 \(K\) 代数：
	\[
		\begin{aligned}
		E_iE_j&=E_jE_i&&\text{当 \(i,j\) 不相邻},\\
		E_i^2E_j-(v+v^{-1})E_iE_jE_i+E_jE_i^2&=0&&\text{当 \(i,j\) 相邻},
		\end{aligned}
	\]
	记为 \(U_v^+(\mathfrak g)\)，称为该参数下量子包络代数的\term[liangzi-zhengbu]{正部}[positive part of a quantum group]。第二组对每个有序相邻对 \((i,j)\) 都要求成立，称为\term[liangzi-Serre-guanxi]{量子 Serre 关系}[quantum Serre relations]。令 \(\deg E_i=\alpha_i\)，得到根格的非负部分上的分次。
\end{definition}
\symidx{\(U_v^+(\mathfrak g)\)：量子包络代数的正部}
这里保留了升算子 \(E_i\)；完整的量子包络代数还含降算子与 Cartan 生成元，\(\mathfrak{sl}_2\) 的例子见\cref{b5:def:B-quantum-sl2}。正部作为代数的上述展示不能单独赋予它此前完整量子群的 Hopf 结构，因为完整余乘法会出现 Cartan 生成元。

\begin{definition}\label{b5:def:B-hall-composition-subalgebra}
	设 \(Q\) 为有限无有向环箭图，\(k=\mathbb F_q\)、\(v=\sqrt q>0\)、\(K=\mathbb Q(v)\)。在扭曲 Hall 代数 \(\mathcal H_v(Q)\) 中，由顶点简单模的类 \([S_i]\) 连同单位元生成的含幺 \(K\) 子代数，称为\term[hecheng-zidaishu]{合成子代数}[composition subalgebra]，记为 \(\mathcal C_v(Q)\)。
\end{definition}
\symidx{\(\mathcal C_v(Q)\)：由顶点简单模生成的 Hall 合成子代数}
每个模虽有合成列，一串简单模的乘积却通常是许多中项的和。仅凭合成列存在，不能把和式中的每个基元素分别取出。因此 \(\mathcal C_v(Q)=\mathcal H_v(Q)\) 是需要证明的结论。

\begin{theorem}[Ringel 定理的 Dynkin 情形]\label{b5:thm:B-ringel-dynkin}
	设 \(Q\) 为一个连通的单边 Dynkin 图的任意定向，\(k=\mathbb F_q\)、\(v=\sqrt q>0\)、\(K=\mathbb Q(v)\)，\(\mathfrak g\) 为相应的复单 Lie 代数。按照\eqref{b5:eq:B-twisted-hall-product}的扭曲及子模、商模次序，赋值
	\[
		\Phi:U_v^+(\mathfrak g)\longrightarrow\mathcal H_v(Q),
		\qquad E_i\longmapsto[S_i]
	\]
	延拓为分次 \(K\) 代数同构。特别地，\(\mathcal C_v(Q)=\mathcal H_v(Q)\)。
\end{theorem}
\begin{proofsketch}
	先证明赋值保持关系。每一条关系只涉及两个顶点；其各项的维数向量在其余顶点上为零，所以参与计数的表示也只支撑在这两个顶点上。不相邻时，上一小节的分裂计数给出交换关系；相邻时，限制后的箭图是 \(1\to2\) 或它的反向，\cref{b5:prop:B-hall-a2-serre}及对偶论证给出两条量子 Serre 关系。因此得到同态 \(\Phi\)，其像按定义就是 \(\mathcal C_v(Q)\)。

	单射性需要比二顶点计算更深的输入。在 Hall 空间上定义对角双线性配对
	\[
		([M],[N])_{\mathrm H}
		=\frac{\delta_{[M],[N]}}{|\operatorname{Aut}_{kQ}(M)|}.
	\]
	Green 的配对相容定理把乘法同短正合列计数给出的余乘法联系起来。通常先在 Hall 代数与量子正部两边都加入记录根格的 Cartan 生成元，使同态扩充为 Hopf 代数同态，再比较两边的 Hopf 配对。配对在简单生成元上的值都是 \(\delta_{ij}/(v^2-1)\)；连同 Cartan 生成元上的配对值及乘法、余乘法相容性，配对的唯一性使 Hall 配对的拉回等于量子群的标准配对。限制回正部后，该配对在这里的非单位根参数 \(v>1\) 上非退化。因此，若 \(a\in\ker\Phi\)，则对每个 \(b\in U_v^+(\mathfrak g)\)，都有 \((a,b)=(\Phi(a),\Phi(b))_{\mathrm H}=0\)，从而 \(a=0\)。这一步所用的配对唯一性及非退化性来自量子群的结构理论，不能仅由 Hall 配对的对角公式推出。

	最后按齐次分量比较维数。记 \(\Delta_+\) 为对应 Dynkin 图的正根集合。有限型表示分类在有限域上仍给每个 \(\beta\in\Delta_+\) 唯一的不可分解模；它可由顶点简单模经反射函子构造，所需核与余核均可在 \(\mathbb F_q\) 上进行。再由 Krull--Schmidt 分解，维数向量为 \(\mathbf d\) 的模由各正根的重数唯一确定，所以
	\[
		\dim_K\mathcal H_v(Q)_{\mathbf d}
		=\#\left\{\,(n_\beta)_{\beta\in\Delta_+}\in\mathbb N^{\Delta_+}
		\mathrel{\Big|}\sum_{\beta\in\Delta_+}n_\beta\beta=\mathbf d\,\right\}.
	\]
	量子 PBW 基定理以每个正根的根向量及其有序幂给出 \(U_v^+(\mathfrak g)\) 的基，故其同次齐次分量具有同一维数。已证的分次单射于是逐次满射，得到同构。概要省略的是 Green 配对的相容证明、量子配对的非退化证明及量子 PBW 基的构造；仅核验量子 Serre 关系只能得到同态，不能替代这些步骤。
\end{proofsketch}
这一定理及所采用的配对证明见\cite[作者公开稿, Section 3.3, Theorem 3.16, Remarks 3.17(i)--(ii)]{Schiffmann2012HallAlgebras}；配对的定义与相容性见同一公开稿的 Section 1.6, Proposition 1.12, Corollary 1.13，量子配对见 Appendix A.4, Theorem A.18。有限型下进一步把不同有限域上的结构常数统一为 Hall 多项式的结果，见\cite[作者公开稿, Section 3.4, Proposition 3.18]{Schiffmann2012HallAlgebras}。这里的同构已经在固定 \(q\) 和 \(v=\sqrt q\) 的条件下成立，不需要先把参数视作不定元。

\ProseParagraphBreak
对单顶点，定理就是 \(K[E]\cong K[u_1]\)，而除幂对应\eqref{b5:eq:B-hall-divided-power}中的带权基元素。对 \(1\to2\)，\(vE_1E_2-E_2E_1\) 的像为 \([P]\)，于是简单根与另一个正根分别落在三个不可分解表示上。一般 Dynkin 情形把这一现象扩展到全部正根，并通过计数乘法记录这些表示之间的扩张。

% * 状态：初稿已完成；新增内容尚未正式审核。


\begin{chaptersummary}
	量子包络代数改变了乘法和余乘法的具体公式，簇变异改变生成元的选择，范畴化则用对象与映射解释数值或环中的等式。三者的联系须通过具体构造来说明。
	\ProseParagraphBreak
	泛型参数的结论可能在单位根处失效，有理变换也可能在特定点无定义。保留工作环、分母、分次和正合关系，才能准确判断一种计算在变化后的环境中是否仍然有效。
\end{chaptersummary}
\BookExercises
\begin{exercise}\label{b5:ex:B:01}
	从量子整数定义证明 \([r+1]_q=q[r]_q+q^{-r}\)，并计算 \([2]_q,[3]_q\)。
\end{exercise}
\begin{solution}
	把两边乘 \(q-q^{-1}\)，右边为 \(q^{r+1}-q^{1-r}+q^{1-r}-q^{-r-1}\)，即所需分子。所以 \([2]_q=q+q^{-1}\)、\([3]_q=q^2+1+q^{-2}\)。这两个表达式已经是 Laurent 多项式，即使原来的分式写法需要排除 \(q=\pm1\)，它们本身仍可特殊化。
\end{solution}

\begin{exercise}\label{b5:ex:B:02}
	写出 \(m=2\) 量子模型的 \(E,F,K\) 矩阵，并指出使正文不可约性证明成立的系数条件。
\end{exercise}
\begin{solution}
	依基 \(v_0,v_1,v_2\)，有 \(K=\operatorname{diag}(q^2,1,q^{-2})\)，\(F\) 在次对角线两个位置为一，\(E\) 在上对角线两个位置均为 \([2]_q\)，其余为零。正文方法要求三个 \(K\) 特征值互异且 \([2]_q\ne0\)。若 \(q\) 不是单位根，就满足这些条件；若 \(q^2=-1\)，则 \([2]_q=0\)，首尾向量的 \(K\) 特征值也相同，但其整数权仍分别为 \(2,-2\)，故原方法失效。
\end{solution}

\begin{exercise}\label{b5:ex:B:03}
	设 \(k\) 为特征零域，\(q\in k^\times\)、\(q^2\ne1\)，\(V\) 为正文中 \(m=1\) 的二维 \(U_q(\mathfrak{sl}_2)\) 模型。在 \(V\otimes_k V\) 中，求一个被 \(E,F\) 零化且被 \(K\) 固定的非零向量。
\end{exercise}
\begin{solution}
	令 \(w=v_0\otimes v_1-q^{-1}v_1\otimes v_0\)。两项的整数权均为 \(1+(-1)=0\)，由 \(\Delta(K)=K\otimes K\)，其 \(K\) 特征值均为 \(q^0=1\)，故 \(Kw=w\)。余乘法给 \(Ew=(1-q^{-1}q)v_0\otimes v_0=0\)。又 \(F(v_0\otimes v_1)=v_1\otimes v_1\)，\(F(v_1\otimes v_0)=qv_1\otimes v_1\)，所以 \(Fw=0\)。因此 \(w\) 张成一个平凡子模。
\end{solution}

\begin{exercise}\label{b5:ex:B:04}
	计算 \(S^2(E),S^2(F),S^2(K)\)，说明对极的平方未必为恒等。
\end{exercise}
\begin{solution}
	反乘法性给 \(S^2(E)=-S(K^{-1})S(E)=KEK^{-1}=q^2E\)，同理 \(S^2(F)=KFK^{-1}=q^{-2}F\)，\(S^2(K)=K\)。故 \(S^2=\operatorname{Int}(K)\)。当 \(q^2\ne1\) 时，从非零二维表示上也能看出它一般不等于恒等。
\end{solution}

\begin{exercise}\label{b5:ex:B:05}
	沿用\cref{b5:prop:B-root-of-unity-filtration}的域、参数及模 \(T\subset W\subset V_\ell\)，记平凡模为 \(L_0\)。计算 \(\operatorname{Hom}_{U_q}(L_0,V_\ell)\) 与 \(\operatorname{Hom}_{U_q}(V_\ell,L_0)\)，并由此证明 \(0\to W\to V_\ell\to L_0\to0\) 也不分裂。
\end{exercise}
\begin{solution}
	从 \(L_0\) 出发的同态由一的像确定，该像必须被 \(E,F\) 零化且被 \(K\) 固定。由于 \(\ker F=kv_\ell=T\)，而 \(Ev_\ell=0\)、\(Kv_\ell=v_\ell\)，全部这样的同态恰为 \(1\mapsto av_\ell\)，\(a\in k\)，故第一个 Hom 空间为一维。
	\ProseParagraphBreak
	同态 \(f:V_\ell\to L_0\) 满足 \(f(v_i)=f(Fv_{i-1})=0\)、\(1\le i\le\ell\)，所以只可能在 \(v_0\) 上非零。反过来，任意这样的线性映射都与 \(K,F\) 作用相容；\(E\) 的像也落在 \(W\) 中，因为 \(Ev_1=0\)，故它同样与 \(E\) 相容。因此第二个 Hom 空间也为一维，由商映射 \(\pi:V_\ell\to V_\ell/W=L_0\) 生成。任何 \(s:L_0\to V_\ell\) 的像都在 \(T\subset W\) 中，于是 \(\pi s=0\)，不能成为恒等映射。这条短正合列没有截面，因而不分裂。
\end{solution}

\begin{exercise}\label{b5:ex:B:06}
	对\cref{b5:prop:B-three-vertex-mutations}的种子先沿位置 \(2,1\) 变异。证明此时交换矩阵为
	\[
		B''=\begin{pmatrix}0&1&-1\\-1&0&0\\1&0&0\end{pmatrix},
	\]
	再沿位置 \(3\) 变异，求第三个位置的新变量和所得矩阵，并解释为什么这一步没有增添顶点 \(1,2\) 之间的箭头。
\end{exercise}
\begin{solution}
	第二次变异反向 \(B'\) 的第一行、列，并把 \(b'_{23}=-1\) 改为 \(-1+[b'_{21}]_+[b'_{13}]_+=0\)，得到所列 \(B''\)。这时第三列只有负条目 \(b''_{13}=-1\)，所以新变量 \(x'_3\) 满足 \(x_3x'_3=1+x'_1\)。代入正文的 \(x'_2=(x_1+x_3)/x_2\) 与 \(x'_1=(x'_2+x_3)/x_1\)，得
	\[
		x'_3=\frac{x_1x_2+x_1+x_3+x_2x_3}{x_1x_2x_3} =\frac{(x_1+x_3)(x_2+1)}{x_1x_2x_3}.
	\]
	这是初始变量的 Laurent 多项式。矩阵第三行、列反号，其余修正项为零，故
	\[
		B'''=\begin{pmatrix}0&1&1\\-1&0&0\\-1&0&0\end{pmatrix}.
	\]
	在 \(B''\) 的箭图中，顶点 \(3\) 是源点，只连向顶点 \(1\)，没有经过 \(3\) 的有向两步路径。因此这一次变异只反向与 \(3\) 相邻的箭头；这是当前箭图的性质，不是任意顶点变异都满足的规则。
\end{solution}

\begin{exercise}\label{b5:ex:B:07}
	设 \(k\) 为特征零域，\(x,y\) 在 \(k\) 上代数独立，\(z_1,\ldots,z_5\in k(x,y)\) 为正文 \(A_2\) 例子的五个簇变量。定义 \(k\) 代数同态
	\[
		\phi:k[Z_1,\ldots,Z_5]\longrightarrow k(x,y),\qquad Z_i\longmapsto z_i.
	\]
	验证恒等式 \(z_1z_3-z_2=1\)、\(z_2z_4-z_3=1\)、\(z_3z_5-z_4=1\)，并说明对应的三条多项式关系是否生成 \(\ker\phi\)。
\end{exercise}
\begin{solution}
	将五个有理式代入，前两式即递推；第三式分子为 \((1+x)(1+y)-(1+x+y)=xy\)，除以 \(xy\) 为一。因此
	\[
		I=(Z_1Z_3-Z_2-1,\ Z_2Z_4-Z_3-1,\ Z_3Z_5-Z_4-1) \subseteq\ker\phi.
	\]
	但另有 \(z_1z_4-z_5=1\)，故 \(Z_1Z_4-Z_5-1\in\ker\phi\)。在点 \((0,-1,0,-1,0)\) 处，\(I\) 的三个生成元均取零，而这个新关系取 \(-1\)，理想 \(I\) 中任一元素都是三个生成元的多项式倍数之和，故在它们的公共零点处必取零；所以这个取值为 \(-1\) 的新关系不属于 \(I\)。因此这三条关系没有生成整个核。
\end{solution}

\begin{exercise}\label{b5:ex:B:08}
	在 \(A_2\) 的初始点 \(x=y=1\) 计算五变量；再在 \(x=y=-1\) 计算，比较哪个相邻簇可作为非零坐标。
\end{exercise}
\begin{solution}
	第一点得到 \(1,1,2,3,2\)，全部非零。第二点得到 \(-1,-1,0,-1,0\)，只有初始簇两变量都非零；其他四个相邻簇至少含一个零。原初始Laurent表达在两点都有意义，四个其他相邻簇为 \(\{z_2,z_3\},\{z_3,z_4\},\{z_4,z_5\},\{z_5,z_1\}\)，每组均含零，不能作为可逆坐标。例如从 \(\{z_2,z_3\}\) 恢复 \(x\) 时需用 \(x=(1+z_2)/z_3\)，第二点的分母为零，尽管原初始坐标 \(x=-1\) 已有确定值。
\end{solution}

\begin{exercise}\label{b5:ex:B:09}
	令 \(z_1=x,z_2=y\)，并取递推 \(z_{i-1}z_{i+1}=1+z_i^2\)。计算 \(z_3,z_4\)，证明这两步仍为 Laurent 多项式。
\end{exercise}
\begin{solution}
	\(z_3=(1+y^2)/x\)，而 \(z_4=\bigl(x^2+(1+y^2)^2\bigr)/(x^2y)\)。分母均为单项式，所以是Laurent表达。这只核验两步，不能据此证明所有变异的Laurent性或有限型。
\end{solution}

\begin{exercise}\label{b5:ex:B:10}
	设 \(A=k[t]/(t^2)\)。比较其有限维模范畴的分裂Grothendieck群与 \(G_0\)，并写出自然满同态的核。
\end{exercise}
\begin{solution}
	不可分解模为 \(k\) 和 \(A\)，故分裂群为 \(\mathbb Z[k]\oplus\mathbb Z[A]\)。合成列 \(0\to k\to A\to k\to0\) 给 \([A]=2[k]\)，所以 \(G_0\cong\mathbb Z[k]\)，自然映射的核由 \([A]-2[k]\) 生成。
\end{solution}

\begin{exercise}\label{b5:ex:B:11}
	在 \(\mathfrak{sl}_2\) 的表示Grothendieck环中，以 \(t=[L_1]\) 表示 \([L_2],[L_3],[L_4]\)，并求 \([L_2]^2\)。
\end{exercise}
\begin{solution}
	递推给 \([L_2]=t^2-1\)、\([L_3]=t^3-2t\)、\([L_4]=t^4-3t^2+1\)。因此 \([L_2]^2=t^4-2t^2+1=[L_4]+[L_2]+[L_0]\)，与完整Clebsch--Gordan分解一致。多项式中的负系数来自换基，并不意味着实际张量分解有负重数。
\end{solution}

\begin{exercise}\label{b5:ex:B:12}
	令 \(C^\bullet\) 为微分内部次数为零的两项分次复形，\(C^0=k\oplus k\langle1\rangle\)、\(C^1=k\langle1\rangle\)，其中 \(k\langle1\rangle\) 表示集中在内部次数一的一维空间。若微分在次数一上为恒等，求其Euler类与上同调分次。
\end{exercise}
\begin{solution}
	Euler类为 \((1+v)-v=1\)。微分把唯一内部次数一部分同构送到后一项，所以 \(H^0\cong k\) 在内部次数零，\(H^1=0\)。若把整个复形移一位，Euler类变成 \(-1\)，而若只把内部次数平移一位，则变成 \(v\)。
\end{solution}
